\bookchapter{الفصل 36: الهاتف والمنزل الذكي — اجعل منزلك ذكيًا}

\chapterintro{يمكن للهاتف أن يفعل شيئًا يزداد أهمية — أن يكون «الجهاز التحكمي» للمنزل بأكمله. من مصباح ومقبس إلى قفل باب، يتحول المنزل الذكي من «لعبة المحترفين» إلى «تجهيز منزلي». يساعدك هذا الفصل على فهم البروتوكولات والمنظومات والمراكز وأطقم البداية، وينتهي بقائمة أمان حتى لا تنقلب عليك «الذكاء».}

\booksection{أولاً: بروتوكولات المنزل الذكي: الاختيار مهم}

قبل شراء الأجهزة الذكية، انظر إلى البروتوكولات المدعومة، لأنها تحدد إمكانية التحكم من هاتفك ومنظومتك. هناك أربعة بروتوكولات رئيسية، ولكل منها خصائصه:

\begin{itemize}
  \item بروتوكول Matter: أصدرته CSA (تحالف معايير الاتصال) في 2022، ويدعمه كبار مثل Apple وGoogle وAmazon وSamsung. هدفه واحد — جعل جميع الأجهزة الذكية تتحدث اللغة نفسها. في الفترة 2023-2026 يتحول بسرعة إلى معيار موحد، ويمكن للأجهزة الداعمة له الاتصال بـ Apple HomeKit وGoogle Home وAmazon Alexa وMi Home في وقت واحد، دون ربطها بمنظومة واحدة.
  \item بروتوكول Thread: «أنبوب» الشبكة الأساسي لـ Matter. منخفض الاستهلاك، وشبكة ذاتية التكوين (إذا انقطع عقدة يعيد الإشارة مسارها)، وسريع الاستجابة دون تأخر. الأجهزة المزودة بموجّه حدود Thread (مثل HomePod mini وApple TV 4K وXiaomi Hub) تمنحك تجربة أكثر استقرارًا.
  \item Wi-Fi: نطاق عريض، واتصال مباشر بالموجّه، لكنه يستهلك طاقة أكثر، وكثرة الأجهزة قد تزدحم القنوات، وهو مناسب للكاميرات والمكنسة الكهربائية التي تنقل بيانات كبيرة.
  \item Bluetooth (BLE): إعداد سهل وموفر للطاقة، لكن المدى قصير ولا يدعم التكوين الذاتي، وهو مناسب للأقفال وأجهزة قياس الحرارة والرطوبة ذات البيانات الصغيرة.
  \item Zigbee: بروتوكول قديم منخفض الاستهلاك بشبكة ذاتية التكوين، يحتاج إلى بوابة مخصصة، والأجهزة رخيصة ومستقرة، لكن التوافق بين العلامات ضعيف، ويتم استبداله تدريجيًا بـ Matter+Thread.
\end{itemize}

\begin{tcbinfo}{المبدأ باختصار}
عند شراء أجهزة ذكية في 2026، فضّل وجود علامة «Matter over Thread» — وهو الخيار الأكثر استدامة حاليًا. أما الأجهزة القديمة التي تدعم Wi-Fi أو Bluetooth فقط فلا داعي للتخلص منها فورًا، بل استخدمها، ثم انتقل إلى Matter عند الشراء الجديد.

\end{tcbinfo}

\booksection{ثانيًا: مقارنة المنظومات الثلاث الكبرى}

البروتوكول هو «اللغة»، والمنظومة هي «الدائرة». حتى لو دعمت جميعها Matter، فأنت لا تزال تختار منظومة رئيسية للأتمتة والتحكم الصوتي. إليك المقارنة:

\begin{center}
\begin{tabularx}{\linewidth}{|X|X|X|X|}
\hline
\textbf{المنظومة} & \textbf{المزايا} & \textbf{العيوب} & \textbf{الفئة المستهدفة} \tabularnewline
\hline
Apple HomeKit & أقوى خصوصية (معالجة محلية)، وتكامل سلس مع iPhone/Siri/الاختصارات & أجهزة أقل وأغلى، وتحتاج إلى HomePod أو Apple TV كمركز & مستخدمو أجهزة Apple في المنزل، المهتمون بالخصوصية والتجربة \tabularnewline
Mi Home (Xiaomi) & أغنى أجهزة وأفضل قيمة، وXiao Ai مفيد & منظومة نسبية مغلقة، وبيانات تُرفع إلى السحابة، وإعلانات أكثر & المبتدئون بميزانية محدودة الراغبون في تجهيز المنزل بالكامل \tabularnewline
Huawei HarmonyOS Connect & تكامل سلس مع هواتف/أجهزة HarmonyOS، وعدد الأجهزة ينمو بسرعة & دعم محدود للأجهزة في الخارج، والمنظومة لا تزال تلاحق & مستخدمو هواتف Huawei، المهتمون بـ«الأجهزة الفائقة» وتعاون الأجهزة \tabularnewline
\hline
\end{tabularx}
\end{center}

\begin{tcbinfo}{توصيات اختيار المنظومة}
اختر المنظومة حسب هاتفك: مستخدمو iPhone يفضلون HomeKit، ومستخدمو Huawei يفضلون HarmonyOS Connect، وباقي مستخدمي Android يجدون Mi Home الأسهل. بفضل بروتوكول Matter، يمكن التحكم في أجهزة المنظومات المختلفة من الهاتف نفسه، دون الحاجة إلى «انحياز» فوري.

\end{tcbinfo}

\booksection{ثالثًا: الهاتف كمركز للمنزل}

في معظم سيناريوهات البداية، الهاتف هو «مركز منزلك» — عبر التطبيقات يمكنك التحكم في الإضاءة والتكييف والستائر والأمان، وبناء سيناريوهات آلية تجعل المنزل «يتحرك من تلقاء نفسه»:

\begin{itemize}
  \item التحكم في الأجهزة: افتح HomeKit/Mi Home/Huawei Smart Life، وشغّل أو أوقف الإضاءة عن بُعد أو محليًا، واضبط درجة حرارة التكييف، وافتح أو أغلق الستائر، وشاهد الكاميرات.
  \item السيناريوهات الآلية: اجمع بين «الشرط الإجراء»، وأشهر ثلاثة سيناريوهات: وضع العودة إلى المنزل: تحديد الموقع + حلول الظلام تشغيل مصباح المدخل + التكييف + إغلاق الستارات + قفل الباب.
  \item وضع المغادرة: آخر شخص يخرج (تحديد الموقع أو قفل الباب) إطفاء الأضواء + إيقاف التكييف + إغلاق المقابس + تشغيل الكاميرات + التفعيل الأمني.
  \item وضع النوم: كل ليلة 23:00 إطفاء الإضاءة الرئيسية + إغلاق الستارات + ضبط التكييف على 26 درجة + تشغيل مصباح الليل.
\end{itemize}

لمس NFC: ضع بطاقة NFC على الباب، ومسها بالهاتف لتفعيل السيناريو (لمسة لفتح ضوء المدخل، ولمسة أخرى لوضع النوم). في iPhone استخدم «الاختصارات»، وفي Android يمكن الكتابة عبر «المحفظة» أو تطبيقات الشركات.

\begin{tcbpractice}{جرّب الآن: أنشئ «وضع العودة إلى المنزل»}
في تطبيق المنزل الذكي على الهاتف، أنشئ سيناريو جديدًا في «الأتمتة»: اختر شرط «الوصول إلى موقع» (المنزل)، وأضف إجراءات «تشغيل مصباح المدخل» و«ضبط التكييف على 26» و«إغلاق الستارات». احفظ، ثم اخرج وجول، وعند عودتك شاهد هل يضيء المصباح من تلقاء نفسه — الإحساس بالنجاح الأولي أفضل من شراء جهاز جديد.

\end{tcbpractice}

\booksection{رابعًا: أطقم البداية الموصى بها}

لا حاجة لتجهيز كل شيء دفعة واحدة، ابدأ بهذه الرباعية لتغطية 80\% من الاحتياجات اليومية:

\begin{itemize}
  \item مكبر صوت ذكي: Xiao Ai (منظومة Mi Home، قيمة ممتازة) / HomePod mini (منظومة Apple، يعمل كمركز Thread) / Alexa (منظومة عالمية). ملاحظة: مكبر الصوت يحدد منظومتك الصوتية، لذا اختر المنظومة أولًا ثم اشترِ السماعة.
  \item مصباح ذكي: Yeelight (Mi Home/Matter، سطوع ودرجة حرارة قابلة للضبط، رخيص) / Philips Hue (ألوان ممتازة ومستقرة، لكنه يحتاج إلى بوابة Hue ويكون أغلى). ملاحظة: تأكد من قاعدة المصباح (E27/B22) وحاجتك لتغيير اللون.
  \item مقبس ذكي: Xiaomi Smart Plug (رخيص، مع إحصاءات الطاقة) / TP-Link Kasa (دعم Matter جيد في الخارج). ملاحظة: عند استخدامه مع الأجهزة عالية القدرة (سخان المياه والمدفأة)، تأكد من القدرة الاسمية ولا تُحمّل المقبس.
  \item قفل باب ذكي: Loock/Xiaomi (قيمة جيدة وتكامل مع Mi Home) / Kaadas (مصنع أقفال تقليدية بهيكل متين). ملاحظة: تأكد من دعم Matter أو منظومتك، وانتبه لعمر البطارية وموضع ثقب المفتاح الطارئ.
\end{itemize}

\begin{tcbinfo}{توصية ترتيب الشراء}
اشترِ مكبر الصوت الذكي أولًا (تحديد المنظومة) ثم المصابيح والمقابس (سهلة الاستخدام وتمنحك تغذية راجعة فورية) وأخيرًا القفل والكاميرا (تتعلقان بالأمان، فاختر بحذر). أضف نوعًا واحدًا في كل مرة، وتوسّع بعد الإتقان، لتجنّب «الشراء دون الاستخدام».

\end{tcbinfo}

\booksection{خامسًا: ملاحظات الأمان}

كلما زادت الأجهزة الذكية في المنزل، زادت سطوح الهجوم. فيما يلي أربع قواعد أساسية، راجعها خطوة بخطوة:

\begin{enumerate}
  \item غيّر كلمة المرور الافتراضية: كلمات المرور الافتراضية للكاميرات والموجّهات والأقفال متاحة للعامة على الإنترنت، فغيّرها فورًا إلى كلمة مرور «قوية وغير متكررة».
  \item حدّث البرنامج الثابت بانتظام: تحديثات الشركات تصلح الثغرات الأمنية، لذا فعّل «التحديث التلقائي» في التطبيق، أو على الأقل تحقق يدويًا شهريًا.
  \item عزل الشبكة: ضع أجهزة IoT على Wi-Fi منفصل (أو شبكة الضيوف)، وافصلها عن الشبكة الرئيسية للهاتف والكمبيوتر. حتى إذا اختُرق جهاز، يصعب تهديد بياناتك المهمة.
  \item حماية الخصوصية: فعّل «مناطق الخصوصية» للكاميرات (مثل عدم تصوير غرفة النوم أو الحمام)، وأوقف التسجيل السحابي 24 ساعة؛ ولا تترك شاشات الكاميرات المنزلية مفتوحة لتطبيقات طرف ثالث لفترة طويلة.
\end{enumerate}

\textbf{لا تقع في هذه الفخاخ}\par

\begin{itemize}
  \item لا تستخدم ميزة «الفتح عن بُعد» للقفل مع الغرباء — فضل كلمة مرور مؤقتة تنتهي فور الاستخدام.
  \item لا توجه الكاميرا مباشرة نحو السرير أو النافذة — فأنت تسرب خصوصيتك أنت.
  \item لا تشترِ أجهزة «بيضاء» غير معروفة من أجل الرخص — غالبًا لا يتم تحديثها، وتُعد بمثابة باب خلفي في منزلك.
\end{itemize}

\begin{tcbsummary}{ملخص الفصل}
\begin{itemize}
  \item Matter+Thread هو الاتجاه الموحد 2023-2026، فعند الشراء الجديد انظر إلى هذه العلامة أولًا
  \item للمنظومات الثلاث مزاياها: Apple تركز على الخصوصية والتجربة، وMi Home على القيمة، وHuawei على تعاون الأجهزة
  \item يمكن استخدام الهاتف كمركز، وبناء سيناريوهات «العودة/المغادرة/النوم» الثلاثة، وسيتحرك المنزل من تلقاء نفسه
  \item أطقم البداية الرباعية: مكبر صوت ذكي ومصباح ومقبس وقفل، وأضفها بالترتيب
  \item غيّر كلمة المرور الافتراضية، وحدّث البرنامج الثابت، واعزل الشبكة، وامنع الكاميرات — هذه القواعد الأربع لا يمكن الاستغناء عنها
\end{itemize}

\end{tcbsummary}
